% ============================================================
% DRAFT NOTE (dataset)
% Contains: The regenerated IFEval sweep, cell definition, coverage, and
% provenance requirements.
% Written now: Settled sweep design and label-provenance rationale.
% Pending: Public artifact location and final manifest identifiers.
% Sources: docs/methodology.md; docs/_why/6_intervention_semantics.md;
% docs/_why/1_measuring_damage.md; docs/_why/3_measuring_quality.md
% ============================================================
\section{Dataset}

The benchmark contains one completed paired intervention for each
$(\text{model}, \text{task}, \text{prompt}, \compressor, \ratio, \switchpos)$
cell. The initial study uses 200 IFEval prompts, three compressors, and four
removal ratios. Its frozen sweep design contains 59,076 switch cells with no
skipped rows. Switch positions are absolute generated-token positions, sampled
on the fixed study grid; their meaning is thus consistent across prompts of
different lengths.

Each row stores the final quality delta and causal features from the reference
stream at the decision boundary. The dataset manifest verifies scorer identity,
ratio semantics, feature timing, exact cell keys, raw-artifact hashes, and
complete prompt coverage. \todo{Add the released manifest version, artifact
location, and exact compressor-to-intervention mapping.}

Historical quality labels were not accepted merely because they shared token
prefixes with the new protocol. Their controls may represent continuous
decoding while their treatments use pressed re-prefill, a relationship that
requires sham and intervention parity. The benchmark therefore uses labels
regenerated under the validated intervention protocol, preserving a direct
causal interpretation for every row.
